The sensing structure is based on the serpentine constriction-resistance sensor developed in \cite{G.O.A.T}, which is taken as a template and adapted to the requirements of the present application. As shown in \cref{fig:serpentine}, the structure consists of a single continuous conductive track, or traxel, printed as a series of closely spaced turns, where the turn diameter $d$ and traxel width $w$ set the spacing between adjacent turns, while the pattern height $h$ and the number of turns $n$ set the overall extent of the serpentine. The over-extrusion that occurs at the turns brings neighbouring turns into contact at those points, as seen in the close-up, so that at rest the turns touch and the current passes directly across the contacts. As the structure is stretched along the sensing direction the contacts progressively open and force the current through the full length of the traxel, which raises the resistance.

\begin{figure}
\centering
\includestandalone[mode=buildmissing]{Standalone/Drawings/sensor_drawing}
    \caption{The serpentine sensing structure pattern is a single continuous traxel of width $w$, wound into $n$ turns of diameter $d$ over a pattern height $h$, with the sensing direction indicated. The close-up shows the contact formed between adjacent turns by the over-extrusion at each turn.}
    \label{fig:serpentine}
\end{figure}

Reproducing such a fine conductive pattern reliably is difficult with a conventional slicer, which converts a solid model into stacked layers and offers only indirect control over the exact path along which material is deposited. The pattern is therefore generated with FullControl, an open-source tool that describes a print directly as explicit G-code by specifying the precise path, extrusion, and speed of the nozzle rather than slicing a model into layers \cite{fullcontrol}. The serpentine is designed as a single defined toolpath, so that the turn diameter and spacing which govern the resting contact are set directly and reproducibly. The fabrication uses a hybrid of the two approaches, in which the TPU substrate and its surface-smoothing ironing pass are produced with a conventional slicer while the FullControl G-code is injected only for the conductive serpentine, combining a uniform substrate with precise deposition of the sensing traxel.

Because the turns are printed close enough to remain in contact, they partially fuse during printing and must be separated before the sensor reaches a usable state, a conditioning step referred to as break-in. After break-in the contacts settle into a stable resting configuration that defines the unloaded resistance of the sensor.

The present application differs from the original use of the sensor mainly in the magnitude of the strains involved, which are up to \SI{800}{\micro\strain} rather than the several-percent range, of order \SIrange{5}{10}{\percent}, for which the template was characterised. At this scale the turns displace only slightly, so usable sensitivity is obtained only if the broken-in resting state already places the contacts close to the point of separation, where a small change in turn spacing produces a large change in the number and area of the contacts and therefore a large relative change in resistance. The same near-separation resting state that yields this sensitivity also leaves the unloaded resistance strongly dependent on fabrication and on the environment, which motivates the baseline, drift, and temperature characterisation that follows and, ultimately, the bridge readout introduced to compensate for it.